% !TEX program = xelatex
% 波动方程与热传导的格林函数
\documentclass[fontset=windows, chinese, 10pt, xcolor=dvipsnames]{beamer}
% Shared preamble for the chapter 10 standalone lectures.
% Each lecture sets \title and \subtitle after this file is read.

\usepackage{amsmath,amssymb,mathtools}
\usepackage{xeCJK}
\usepackage{zhnumber}
\usepackage{tikz}
\usepackage{tikz-3dplot}
\usetikzlibrary{arrows.meta,calc,positioning}

\usetheme{metropolis}
\usepackage{beamercolorthemeHistNanjing}

\newcommand{\red}[1]{\textcolor{red}{#1}}
\newcommand{\blue}[1]{\textcolor{blue}{#1}}
\newcommand{\green}[1]{\textcolor{green}{#1}}
\newcommand{\gray}[1]{\textcolor{gray}{#1}}
\newcommand{\orange}[1]{\textcolor{orange}{#1}}

\newcommand{\sssub}[1]{ {\mathrm{#1} }}  % script size subscript\\\
\newcommand{\funcd}[2]{ \frac{d{#1}}{d{#2}}} %function full derivative
\newcommand{\funcpd}[2]{ \frac{\partial{#1}}{\partial{#2}}} %function partial derivative
\newcommand{\fnald}[2] { \frac{\delta {#1}}{\delta {#2}}} % functional derivative

\renewcommand{\Re}{\operatorname{Re}}
\renewcommand{\Im}{\operatorname{Im}}
\def\ext{\sssub{ext}}
\def\x{\sssub{x}}
\def\xc{\sssub{xc}}
\def\H{\sssub{H}}
\def\Hx{\sssub{Hx}}
\def\Hxc{\sssub{Hxc}}
\def\c{\sssub{c}}
\def\s{\sssub{s}}

% \def\ext{\sssub{ext}}
% \def\ext{\sssub{ext}}
% \def\ext{\sssub{ext}}
% \def\ext{\sssub{ext}}
% \defsssub{ext}


% definitions of bold faced symbols
\def\brpppp{{\mathbf{r}^{\prime\prime\prime\prime}}}
\def\brppp{{\mathbf{r}^{\prime\prime\prime}}}
\def\brpp{{\mathbf{r}^{\prime\prime}}}
\def\brp{{\mathbf{r}^{\prime}}}
\def\bzp{{\mathbf{z}^{\prime}}}
\def\bxp{{\mathbf{x}^{\prime}}}
\def\tp{{{t}^{\prime}}}
\def\tpp{{{t}^{\prime\prime}}}
\def\tppp{{{t}^{\prime\prime\prime}}}

\def\tbr{{\tilde{\mathbf{r}}}}
\def\bk{{\mathbf{k}}}
\def\br{{\mathbf{r}}}
\def\bp{{\mathbf{p}}}
\def\bq{{\mathbf{q}}}
\def\bv{{\mathbf{v}}}
\def\bz{{\mathbf{z}}}
\def\mbf{{\mathbf{f}}}

\def\bx{{\mathbf{x}}}
\def\bR{{\mathbf{R}}}
\def\bM{{\mathbf{M}}}
\def\bP{{\mathbf{P}}}
\def\bT{{\mathbf{T}}}
\def\bK{{\mathbf{K}}}
\def\bA{{\mathbf{A}}}
\def\bB{{\mathbf{B}}}
\def\bD{{\mathbf{D}}}
\def\bE{{\mathbf{E}}}
\def\bF{{\mathbf{F}}}
\def\bG{{\mathbf{G}}}
\def\bH{{\mathbf{H}}}
\def\bI{{\mathbf{I}}}
\def\bJ{{\mathbf{J}}}
\def\bP{{\mathbf{P}}}
\def\bS{{\mathbf{S}}}
\def\bV{{\mathbf{V}}}
\def\bX{{\mathbf{X}}}
\def\bY{{\mathbf{Y}}}
\def\dd{{\mathop{}\!\mathrm{d}}}
\def\dr{{\!\!\dd\br\,}}
\def\dx{{\!\!\dd\bx\,}}
\def\dy{{\!\!\dd\by\,}}
\def\dk{{\!\!\dd\bk\,}}
\def\dG{{\!\!\dd\bG\,}}


\def\half{{\frac{1}{2}}}
\def\tr{{\mathrm{tr}}}


\def\Arg{{\mathrm{Arg}\;}}
\def\arg{{\mathrm{arg}\;}}
\def\Res{{\mathrm{Res}\;}}

% unit-vector basis (bold under XeLaTeX; \mathbf can look non-bold)
\newcommand{\ue}{\boldsymbol{e}}
\newcommand{\un}{\hat{\boldsymbol{n}}}


\author{罗凯}
\institute{南京理工大学物理学院}
\date{\zhdigits{\the\year},\zhganzhinian{\the\year}年秋}

\usebackgroundtemplate{%
\begin{tikzpicture}[remember picture,overlay]
  \node[opacity=0.05, at=(current page.center)] {
    \includegraphics[width=0.6\paperwidth]{njust-logo.jpeg}
  };
\end{tikzpicture}}


\title{格林函数法}
\subtitle{第三讲：波动与热传导}

\begin{document}

\begin{frame}
  \titlepage
\end{frame}

\begin{frame}{本讲提纲}
  含时问题的格林函数是单位冲量的响应。因果性要求响应出现在冲量之后。
  \begin{enumerate}
    \item 波动方程的表示公式，初值单独成一项；
    \item 无限长弦上，光锥里 $G=1/(2a)$；
    \item 驱动力频率碰上固有频率时出现共振；
    \item 热核是高斯函数，半直线用奇延拓或偶延拓。
  \end{enumerate}
\end{frame}

\section{波动}

\begin{frame}{单位冲量}
  \[
    u_{tt}-a^2\nabla^2 u=f(\mathbf{r},t),
  \]
  初值 $u|_{t=0}=\varphi$，$u_t|_{t=0}=\psi$。连续力是一串冲量。格林函数满足
  \[
    G_{tt}-a^2\nabla^2 G
    =\delta(\mathbf{r}-\mathbf{r}_0)\,\delta(t-t_0),
  \]
  边界条件齐次，并且 $G|_{t=0}=G_t|_{t=0}=0$。$t<t_0$ 时 $G=0$。
\end{frame}

\begin{frame}{三项}
  原方程和 $G$ 的方程相减，在 $T\times[0,t]$ 上积分。时间部分是全导数，空间部分用格林第二恒等式，
  \begin{align*}
    u(\mathbf{r},t)
    &=\iiint\int_0^{t} G\,f\,\mathrm{d}V_0\,\mathrm{d}t_0\\
    &\quad+a^2\iint\int_0^{t}
        \bigl(G\partial_{n_0}u-u\partial_{n_0}G\bigr)
        \,\mathrm{d}S_0\,\mathrm{d}t_0\\
    &\quad+\iiint\bigl[G u_{t_0}-u G_{t_0}\bigr]_{t_0=0}\,\mathrm{d}V_0.
  \end{align*}
  三项依次是外力、边界数据和初值。$G$ 的边界条件与问题同一类型，未知的边界组合不再出现。
\end{frame}

\begin{frame}{时间反演}
  \[
    G(\mathbf{r},t;\mathbf{r}_0,t_0)
    =G(\mathbf{r}_0,-t_0;\mathbf{r},-t).
  \]
  波动方程对 $t\to -t$ 不变。因果性把未来的源变成过去的响应，所以互易性要带着时间反演。
\end{frame}

\section{弦}

\begin{frame}{冲量换成初速度}
  对波动方程，零初位移和单位冲量意味着
  \[
    G\big|_{t=\tau+0}=0,
    \qquad
    G_t\big|_{t=\tau+0}=\delta(x-\xi).
  \]
  热方程对时间是一阶的，跳跃发生在函数值本身，$G|_{t=\tau+0}=\delta(x-\xi)$。
\end{frame}

\begin{frame}{无限长弦}
  达朗贝尔公式把时间原点挪到 $\tau$。初位移为零、初速度为 $\delta(x-\xi)$ 时，依赖区间里的积分是 $1/(2a)$，区间外是零，
  \begin{equation*}
    \boxed{\displaystyle
    G=
    \begin{cases}
      1/(2a), & |x-\xi|<a(t-\tau),\\
      0, & |x-\xi|>a(t-\tau).
    \end{cases}}
  \end{equation*}
  扰动不超过速度 $a$。
\end{frame}

\begin{frame}{叠加}
  \[
    u(x,t)
    =\int_0^{t}\int_{x-a(t-\tau)}^{x+a(t-\tau)}
      \frac{1}{2a}\,f(\xi,\tau)\,\mathrm{d}\xi\,\mathrm{d}\tau.
  \]
  这就是齐次化原理里的那个积分。每个冲量只影响它自己的光锥。
\end{frame}

\begin{frame}{共振}
  弦长 $\ell$，两端自由，外力 $A\cos(\pi x/\ell)\sin\omega t$，初值为零。外力已经是 $n=1$ 的余弦模态。令 $u=q(t)\cos(\pi x/\ell)$，$\omega_0=\pi a/\ell$，
  \[
    q''+\omega_0^2 q=A\sin\omega t,
    \qquad q(0)=q'(0)=0.
  \]
  $\omega\neq\omega_0$ 时
  \[
    q=\frac{A}{\omega_0^2-\omega^2}
      \Bigl(\sin\omega t-\frac{\omega}{\omega_0}\sin\omega_0 t\Bigr).
  \]
  $\omega=\omega_0$ 时
  \[
    q=\frac{A}{2\omega_0^2}
      \bigl(\sin\omega_0 t-\omega_0 t\cos\omega_0 t\Bigr).
  \]
  振幅随 $t$ 线性增长。
\end{frame}

\section{热传导}

\begin{frame}{热方程的表示}
  \[
    u_t-a^2\nabla^2 u=f,
    \qquad u|_{t=0}=\varphi.
  \]
  热方程对时间是一阶的，表示公式只留下一个初值项，
  \begin{align*}
    u
    &=\iiint\int_0^{t} G f\,\mathrm{d}V_0\,\mathrm{d}t_0\\
    &\quad+a^2\iint\int_0^{t}
        \bigl(G\partial_{n_0}u-u\partial_{n_0}G\bigr)
        \,\mathrm{d}S_0\,\mathrm{d}t_0\\
    &\quad+\iiint\bigl[u G\bigr]_{t_0=0}\,\mathrm{d}V_0.
  \end{align*}
\end{frame}

\begin{frame}{无限长杆的热核}
  初值是 $\delta(x-\xi)$ 时，
  \begin{equation*}
    \boxed{\displaystyle
    G
    =\frac{1}{2a\sqrt{\pi(t-\tau)}}
      \exp\Bigl(-\frac{(x-\xi)^2}{4a^2(t-\tau)}\Bigr),
    \quad t>\tau.}
  \end{equation*}
  只要 $t>\tau$，$G$ 在整条直线上都不为零。按这个方程，信号的传播速度是无限的。
\end{frame}

\begin{frame}{半直线}
  $x>0$、端点 $u|_{x=0}=0$ 时，把热核相减，
  \[
    G\propto
    \exp\Bigl(-\frac{(x-\xi)^2}{4a^2(t-\tau)}\Bigr)
    -\exp\Bigl(-\frac{(x+\xi)^2}{4a^2(t-\tau)}\Bigr).
  \]
  $x=0$ 时两项相同，$G=0$。端点绝热 $u_x|_{x=0}=0$ 时改成相加，偶函数的导数在端点为零。
\end{frame}

\begin{frame}{有限杆}
  $0<x<\ell$、两端为零时，把 $\delta(x-\xi)$ 展成正弦级数，每个模态再乘衰减因子，
  \[
    G
    =\frac{2}{\ell}\sum_{n=1}^{\infty}
      \sin\frac{n\pi x}{\ell}
      \sin\frac{n\pi\xi}{\ell}
      \exp\Bigl(-a^2\bigl(\tfrac{n\pi}{\ell}\bigr)^2(t-\tau)\Bigr).
  \]
  $t\to\tau^+$ 时它回到 $\delta(x-\xi)$。这和分离变量写出的解是同一件事。
\end{frame}

\begin{frame}{这一讲留下的公式}
  波动格林函数在光锥里是 $1/(2a)$，锥外是零。互易性要带着 $t\to -t$。
  \vspace{0.35em}
  热核是高斯函数。端点为零用奇延拓，端点绝热用偶延拓。有限杆上，每个正弦模态乘 $e^{-a^2 k_n^2(t-\tau)}$。
\end{frame}

\end{document}
